% Auto-generated from Chapter-10-Corvin-Glimpsed-at-Border.md
% Source: C:\Users\PCGAME\Desktop\story&universes\chain-weapon-story\finished-manuscript\manuscriptbaseforchapters\Chapter-10-Corvin-Glimpsed-at-Border.md

\chapter{Corvin Glimpsed at Border}
\renewcommand{\CurrentChapterNum}{10}
\POV{\glyphAisen}{Aisen}

The border town of Gresik sat on the map like a bruise—a discolored sprawl of timber palisades and mud-brick warehouses where three trade roads met and none of them wanted to stay. It was the kind of place that existed because geography demanded it, not because anyone had chosen to build something beautiful there, and it wore that reluctance in every plank and nail. The air tasted of livestock and the sulfurous tang of the hot springs that fed the public baths, and beneath that, a mineral sharpness that Aisen had learned to associate with altitude—the mountains were closer here, pressing their cold weight against the sky like something watching from a height.

He was thirteen, nearly fourteen, and he had come with Harald on what his father called a supply run, though Aisen suspected the word concealed a more complicated errand. The wagon carried the usual provisions—salt, lamp oil, the coarse wheat flour that baked into bread dense enough to serve as currency in the lean months—but Harald had also packed two sealed letters in his coat pocket, and he had asked Aisen to wear his good boots, the ones with the reinforced soles. These were the kinds of details that did not survive in isolation. Paired together, they meant that someone in Gresik expected them, and that Harald anticipated the possibility of leaving quickly.

The town's market was louder than Valdimere's by an order of magnitude—not busier, precisely, but unregulated in a way that gave each voice permission to compete. Drovers shouted prices over the bleating of goats penned in wooden crates. A farrier's hammer rang against an anvil in a rhythm that syncopated against the clatter of wagon wheels on cobblestones worn smooth by decades of traffic going nowhere it wanted to go. Aisen catalogued these sounds the way he always catalogued them: ingredients in a recipe he was still learning to cook. The noise told him things. A market that shouted was a market where people did not trust the quality of their goods to speak for itself.

Harald moved through the crowd with the particular alertness of a man who had learned to be comfortable in uncomfortable places. He delivered his letters—one to a dry goods merchant with a limp, one to a woman who ran a cartwright's shop and spoke to Harald in a dialect Aisen could not follow—and then steered them toward the northern edge of town, where the palisade gave way to open ground and the road climbed into terrain too rough for wagons.

"We'll eat here," Harald said, nodding toward a wayhouse that squatted near the gate like a toad waiting for flies. "Stay close to me today."

It was the second sentence that sharpened Aisen's attention. Harald said \emph{stay close} the way other fathers said \emph{I love you}—rarely, and only when the weight of it could not be carried by silence alone.